\documentclass{beamer}
\usepackage{tikz,amsmath,amssymb,hyperref,graphicx,stackrel,animate,tipa}
\usetikzlibrary{positioning,shadows,arrows,shapes,calc,dsp,chains}
\newcommand{\argmax}{\operatornamewithlimits{argmax}}
\newcommand{\argmin}{\operatornamewithlimits{argmin}}
\mode<presentation>{\usetheme{Frankfurt}}
\AtBeginSection[]
{
  \begin{frame}<beamer>
    \frametitle{Outline}
    \tableofcontents[currentsection,currentsubsection]
  \end{frame}
}
\title{Lecture 27: Recursive Least Squares}
\author{Mark Hasegawa-Johnson}
\date{ECE 401: Signal Processing}  
\begin{document}

% Title
\begin{frame}
  \maketitle
\end{frame}

% Title
\begin{frame}
  \tableofcontents
\end{frame}


%%%%%%%%%%%%%%%%%%%%%%%%%%%%%%%%%%%%%%%%%%%%%%%%%%%%%%%%%%%%%%%%%%%%%
\section[LPC]{Review: LPC}

\begin{frame}
  \frametitle{Linear predictive coding (LPC): Synthesis model}

  A linearly predictable signal is one whose samples, $s[n]$, can be estimated
  as the linear prediction of $p$ previous samples plus some small error term $e[n]$:
  \begin{align*}
    s[n] = e[n]+\sum_{k=1}^p a_k s[n-k]
  \end{align*}
\end{frame}

\begin{frame}
  \frametitle{LPC: The Solution!}

  We will choose $a_k$ to minimize the mean squared error (MSE):
  \begin{align*}
    \mathcal{E} &= \sum_m e^2[m]\\
    &= \sum_m \left(s[m]-\sum_{k=1}^p a_ks[m-k]\right)^2
  \end{align*}
  If you differentiate w.r.t. $a_i$, you get $p$ linear equations in
  $p$ unknowns:
  \begin{align*}
    r_s[i] &= \sum_{k=1}^p a_kr_s[k-i]\\
    &\text{\ldots where\ldots}\\
    r_s[i] &\equiv \sum_m s[m-i]s[m] 
  \end{align*}

\end{frame}

\begin{frame}
  \frametitle{The LPC solution as a matrix equation}

  \begin{align*}
    \boldsymbol{\gamma} &= \mathbf{R}\mathbf{a}
  \end{align*}
  \begin{align*}
    \boldsymbol{\gamma}=\left[\begin{array}{c}r_s[1]\\\vdots\\r_s[p]\end{array}\right],~~~
    \mathbf{R}=\left[\begin{array}{ccc}r_s[0]&\cdots&r_s[p-1]\\
        \vdots&\ddots&\vdots\\r_s[p-1]&\cdots&r_s[0]\end{array}\right],~~~
    \mathbf{a}=\left[\begin{array}{c}a_1\\\vdots\\a_p\end{array}\right]
  \end{align*}
\end{frame}
    
%%%%%%%%%%%%%%%%%%%%%%%%%%%%%%%%%%%%%%%%%%%%%%%%%%%%%%%%%%%%%%%%%%%%%
\section[RLS]{Recursive Least Squares}

\begin{frame}
  \frametitle{What if the autocorrelation changes rapidly?}

  \begin{itemize}
  \item LPC assumes that the autocorrelation and LPC coefficients are
    stationary for at least one frame.
    \begin{align*}
      s[n] = e[n]+\sum_{k=1}^p a_k s[n-k]
    \end{align*}
  \item What if the autocorrelation is changing quickly, so that we
    need to update the coefficients every sample?
    \begin{align*}
      s[n] = e[n]+\sum_{k=1}^p a_k[n] s[n-k]
    \end{align*}
  \end{itemize}
\end{frame}

\begin{frame}
  \frametitle{Recursive least squares (RLS)}

  \begin{itemize}
  \item RLS, independently invented by Gauss (1850s) and Lanczos
    (1952), recomputes the autocorrelation $r_s[i,k;n]$ and predictor
    coefficients $a_k[n]$ for every sample!
  \item The word ``recursive'' describes our estimate of the
    autocorrelation:
    \begin{align*}
      r_s[i,k;n] &= s[n-i]s[n-k] + \lambda r_s[i,k;n]
    \end{align*}
    \ldots where $\lambda<1$ is the forgetting factor.
  \end{itemize}
\end{frame}

\begin{frame}
  \frametitle{The forgetting factor, $\lambda$}

  The forgetting factor, $\lambda$, gives a weight of $\lambda^0=1.0$
  to the current sample, and progressively decreasing weights to older
  samples.  For example, here's the curve for $\lambda=0.9$:
  \centerline{\includegraphics[width=0.8\textwidth]{Geometrique_exponentielle.png}}
  \centerline{\tiny{\url{https://commons.wikimedia.org/wiki/File:Geometrique_exponentielle.png}}}
\end{frame}
    
\begin{frame}
  \frametitle{The RLS error term}

  \begin{itemize}
  \item The error we want to minimize is
    $e^2[n]=(s[n]-\sum_ka_k[n]s[n-k])^2$.
  \item We could set that to zero with $a_1[n]=s[n]/s[n-1]$!  It's too easy:
    it has only learned the current sample, it hasn't learned anything
    about the usystem that produced the samples
  \item A more useful definition: pretend that $a_k[n]$ has been
    applied for all recent samples, and minimize a weighted error
    term:
    \begin{align*}
      \mathcal{E} &= \sum_{m=0}^\infty \lambda^me^2[n-m]\\
      &= \sum_{m=0}^\infty\lambda^m\left(s[n-m]-\sum_ka_k[n]s[n-m-k]\right)^2
    \end{align*}
  \end{itemize}
\end{frame}

\begin{frame}
  \frametitle{Adaptive filtering}

  RLS is a type of adaptive filter, in which the desired signal,
  $d_k=s[n]$, is predicted from the past samples
  $\mathbf{x}_k=[s[n-1],\ldots,s[n-p]]^T$:

  \centerline{\includegraphics[width=\textwidth]{Adaptive_Filter_Compact.png}}
  \centerline{\tiny{\url{https://commons.wikimedia.org/wiki/File:Adaptive_Filter_Compact.png}}}
\end{frame}


\begin{frame}
  \frametitle{Differentiating the RLS error term}

  Differentiating $\mathcal{E}$ w.r.t. $a_i$ gives us
  \begin{align*}
    \frac{\partial\mathcal{E}}{\partial{a}_i}  
    &= r_s[0,i;n]-\sum_ka_k[n]r_s[i,k;n]
  \end{align*}
  \ldots and setting it to zero gives us\ldots
  \begin{align*}
    r_s[0,i;n] &= \sum_ka_k[n]r_s[i,k;n]
  \end{align*}
  \ldots where \ldots
  \begin{align*}
    r_s[i,k;n] &\equiv s[n-i]s[n-k] + \lambda r_s[i,k;n-1]\\
    &= \sum_{m=0}^\infty \lambda^m s[n-m-i]s[n-m-k]
  \end{align*}
\end{frame}

\begin{frame}
  \frametitle{RLS: The solution!}

  The solution that minimizes the RLS error is $r_s[0,i;n]=\sum_ka_k[n]r_s[i,k;n]$, i.e.,
  \begin{align*}
    \boldsymbol{\gamma}[n] &= \mathbf{R}[n]\mathbf{a}[n]
  \end{align*}
  \ldots where \ldots
  \begin{align*}
    \boldsymbol{\gamma}[n]=\left[\begin{array}{c}r_s[0,1;n]\\\vdots\\r_s[0,p;n]
      \end{array}\right],~~~
    \mathbf{R}[n]=\left[\begin{array}{ccc}r_s[1,1;n]&\cdots&r_s[1,p-1;n]\\
        \vdots&\ddots&\vdots\\r_s[p-1,1;n]&\cdots&r_s[p,p;n]\end{array}\right],~~~
    \mathbf{a}[n]=\left[\begin{array}{c}a_1[n]\\\vdots\\a_p[n]\end{array}\right]
  \end{align*}
\end{frame}

%%%%%%%%%%%%%%%%%%%%%%%%%%%%%%%%%%%%%%%%%%%%%%%%%%%%%%%%%%%%%%%%%%%%%
\section[Matrices]{Matrix recursion and vector prediction}

\begin{frame}
  \frametitle{Matrix recursion}

  Each autocorrelation is recursive:
  \begin{align*}
    r_s[i,k;n] = s[n-i]s[n-k] + \lambda r_s[i,k;n-1]
  \end{align*}
  \ldots and therefore $\mathbf{R}[n]$ and $\boldsymbol{\gamma}[n]$
  are also recursive:
  \begin{align*}
    \mathbf{R}[n]&=\mathbf{s}[n]\mathbf{s}^T[n] + \lambda\mathbf{R}[n-1]\\
    \boldsymbol{\gamma}[n] &= \mathbf{s}[n]s[n] + \lambda\boldsymbol{\gamma}[n-1]
  \end{align*}
  \ldots where $\mathbf{s}[n]$ is the delayed-signal vector
  \begin{align*}
    \mathbf{s}[n] = \left[\begin{array}{c}s[n-1]\\\vdots\\s[n-p]\end{array}\right]
  \end{align*}
\end{frame}

\begin{frame}
  \frametitle{Vector prediction}

  Notice that, with these definitions, the predicted value of $s[n]$ is
  \begin{align*}
    \tilde{s}[n] &= \sum_k a_k[n] s[n-k]\\
    &= \mathbf{s}^T[n]\mathbf{a}[n]
  \end{align*}
  The prediction error is
  \begin{align*}
    e[n] &= s[n]-\mathbf{s}^T[n]\mathbf{a}[n]
  \end{align*}
\end{frame}
  
\begin{frame}
  \frametitle{RLS problem summary}

  To minimize the forgetting-weighted memory
  \begin{align*}
    \mathcal{E} &= \sum_m\lambda^m\left(s[n-m]-\mathbf{s}^T[n-m]\mathbf{a}[n]\right)^2
  \end{align*}
  \ldots we find coefficients $\mathbf{a}[n]$ according to
  \begin{align*}
    \mathbf{a}[n] &= \mathbf{R}^{-1}[n]\boldsymbol{\gamma}[n]
  \end{align*}
  \ldots where $\mathbf{R}[n]$ and $\boldsymbol{\gamma}[n]$ are
  defined recursively from $s[n]$ and
  $\mathbf{s}[n]=[s[n-1],\ldots,s[n-p]]^T$
  \begin{align*}
    \mathbf{R}[n]&=\mathbf{s}[n]\mathbf{s}^T[n] + \lambda\mathbf{R}[n-1]\\
    \boldsymbol{\gamma}[n] &= \mathbf{s}[n]s[n] + \lambda\boldsymbol{\gamma}[n-1]
  \end{align*}
\end{frame}

\begin{frame}
  \frametitle{Try the quiz!}

  Go to the course web page and try the quiz!
\end{frame}

%%%%%%%%%%%%%%%%%%%%%%%%%%%%%%%%%%%%%%%%%%%%%%%%%%%%%%%%%%%%%%%%%%%%%%%%%%%%%%%%
\section{How to make RLS efficient and stable}

\begin{frame}
  \frametitle{Is RLS efficient and stable?}

  RLS has two big problems:
  \begin{itemize}
  \item \textbf{Inefficiency:} We need to invert $\mathbf{R}[n]$ at
    every sample time
  \item \textbf{Instability:} Sometimes $\mathbf{R}[n]$ might be
    nearly singular, in which case $\mathbf{R}^{-1}[n]$ and
    $\mathbf{a}[n]$ might become nearly infinite
  \end{itemize}
  Both problems are solved in the same way, using the
  \textbf{stable RLS algorithm}.
\end{frame}

\begin{frame}
  \frametitle{The stable RLS algorithm}

  At every sample $n$:
  \begin{enumerate}
  \item Compute the \textit{a priori} error $\epsilon[n]$:
    \begin{align*}
      \epsilon[n] = s[n] - \mathbf{s}^T[n]\mathbf{a}[n-1]
    \end{align*}
  \item Compute the update gain vector $\mathbf{g}[n]$
  \item Update the coefficients
    \begin{align*}
      \mathbf{a}[n] = \mathbf{a}[n-1]+\mathbf{g}[n]\epsilon[n]
    \end{align*}
  \end{enumerate}
\end{frame}

\begin{frame}
  \frametitle{The gain vector and the precision matrix}

  \begin{itemize}
  \item The gain vector is defined as
    $\mathbf{g}[n]=\mathbf{R}^{-1}[n]\mathbf{s}[n]$, which is not very
    useful because it still depends on $\mathbf{R}^{-1}[n]$.  With the
    help of some linear algebra, though, we can rewrite it as
    \begin{align*}
      \mathbf{g}[n] =
      \frac{\mathbf{P}[n-1]\mathbf{s}[n]}{\lambda+\mathbf{s}^T[n]\mathbf{P}[n-1]\mathbf{s}[n]}
    \end{align*}
  \item To make the recursion work, we then need to update the
    precision matrix.  The precision matrix is defined as
    $\mathbf{P}[n]=\mathbf{R}^{-1}[n]$, which, again, is not very
    useful, but using some linear algebra we can rewrite it as
    \begin{align*}
      \mathbf{P}[n]=\frac{1}{\lambda}\left(\mathbf{P}[n-1]
      -\mathbf{g}[n]\mathbf{s}^T[n]\mathbf{P}[n-1]\right)
    \end{align*}
    \item Notice that both of these formulas are computationally
      efficient (the worst computation, the matrix-vector
      multiplication, is $\mathcal{O}\left\{p^2\right\}$) and, as long
      as $\lambda>0$, both are stable (finite inputs always give
      finite outputs).
  \end{itemize}
\end{frame}

\begin{frame}
  \textbf{Simulation example}

  \begin{itemize}
  \item Stationary for 500 samples, then sudden change
  \item Wiener filter (black) knows the true $r_s[n]$
  \item LMS (least-mean-squares), SG (stochastic gradient), and
    RLS estimate $r_s[n]$ from the sample values
  \item $\lambda=1.0$ converges to $\mathcal{E}=0$ for the first
    system, but $\lambda=0.8$ adapts better when the system changes.
  \end{itemize}
  \centerline{\includegraphics[width=0.6\textwidth]{SG_RLS_LMS_chan_var.png}}
  \centerline{\tiny{\url{https://commons.wikimedia.org/wiki/File:SG_RLS_LMS_chan_var.png}}}
\end{frame}
        

%%%%%%%%%%%%%%%%%%%%%%%%%%%%%%%%%%%%%%%%%%%%%%%%%%%%%%%%%%%%%%%%%%%5
\section{Summary}

\begin{frame}
  \frametitle{RLS problem summary}

  To minimize the forgetting-weighted memory
  \begin{align*}
    \mathcal{E} &= \sum_m\lambda^m\left(s[n-m]-\mathbf{s}^T[n-m]\mathbf{a}[n]\right)^2
  \end{align*}
  \ldots we find coefficients $\mathbf{a}[n]$ according to
  \begin{align*}
    \mathbf{a}[n] &= \mathbf{R}^{-1}[n]\boldsymbol{\gamma}[n]
  \end{align*}
  \ldots where $\mathbf{R}[n]$ and $\boldsymbol{\gamma}[n]$ are
  defined recursively from $s[n]$ and
  $\mathbf{s}[n]=[s[n-1],\ldots,s[n-p]]^T$
  \begin{align*}
    \mathbf{R}[n]&=\mathbf{s}[n]\mathbf{s}^T[n] + \lambda\mathbf{R}[n-1]\\
    \boldsymbol{\gamma}[n] &= \mathbf{s}[n]s[n] + \lambda\boldsymbol{\gamma}[n-1]
  \end{align*}
\end{frame}

\begin{frame}
  \frametitle{The stable RLS algorithm}

  At every sample $n$:
  \begin{enumerate}
  \item Compute the \textit{a priori} error $\epsilon[n]$:
    \begin{align*}
      \epsilon[n] = s[n] - \mathbf{s}^T[n]\mathbf{a}[n-1]
    \end{align*}
  \item Compute the gain and precision as
    \begin{align*}
      \mathbf{g}[n] &=
      \frac{\mathbf{P}[n-1]\mathbf{s}[n]}{\lambda+\mathbf{s}^T[n]\mathbf{P}[n-1]\mathbf{s}[n]}\\
      \mathbf{P}[n] &=\frac{1}{\lambda}\left(\mathbf{P}[n-1]
      -\mathbf{g}[n]\mathbf{s}^T[n]\mathbf{P}[n-1]\right)
    \end{align*}
  \item Update the coefficients
    \begin{align*}
      \mathbf{a}[n] = \mathbf{a}[n-1]+\mathbf{g}[n]\epsilon[n]
    \end{align*}
  \end{enumerate}
\end{frame}



\end{document}

